\documentclass[7pt,landscape]{article}
\usepackage[utf8]{inputenc}
\usepackage[margin=0.18in]{geometry}
\usepackage{multicol}
\usepackage{titlesec}
\usepackage{enumitem}
\usepackage{xcolor}

\definecolor{navy}{RGB}{0, 51, 102}
\definecolor{linuxc}{RGB}{160, 88, 24}
\definecolor{winc}{RGB}{45, 95, 140}
\setlength{\columnsep}{0.25in}
\setlist{nosep, leftmargin=10pt}
\titlespacing*{\section}{0pt}{1pt}{0.5pt}
\setlength{\parindent}{0pt}
\linespread{0.90}

\titleformat{\section}{\small\bfseries\color{navy}}{\thesection}{0.5em}{}[\titlerule]
\pagestyle{empty}

\begin{document}
\textbf{Role Reference: SIEM / Enumeration Lead} \hfill \textit{Last Updated: Aug 2026}
\raggedcolumns
\begin{multicols*}{3}

\section{ROLE RESPONSIBILITIES}
Standing up and tuning the SIEM. Calling out alerts. Confirming scope across every host, not just the one under attack.

\section{ORDER OF OPERATIONS}
\textbf{Round start:} Confirm Security Onion (and Wazuh, if deployed) is up, sensors capturing, agents reporting. \\
\textbf{Tuning pass:} Reduce known-noisy alerts before the round gets busy. \\
\textbf{Ongoing:} Work the alert queue continuously. \\
\textbf{On a reported foothold:} Search every host for the same account, IOC, or timestamp. Report scope to the Captain. \\
\textbf{After eviction:} Keep watching for it to return.

\section{SECURITY ONION: SETUP CHECKS}
\begin{itemize}
    \item \texttt{so-status}: every component running.
    \item \texttt{so-allow}: whitelist your IP if the web UI refuses connection.
    \item Monitor interface (not management) is the one seeing traffic; packet counts climbing.
    \item \texttt{so-user}: manage additional analyst accounts on the box itself.
\end{itemize}

\section{HUNT: KEY QUERIES}
\begin{itemize}
    \item \textit{Lateral movement:} \texttt{event.dataset:conn AND source.ip:"[host]"} \newline \texttt{AND destination.ip:10.0.0.0/8}
    \item \textit{Exfiltration:} \texttt{event.dataset:conn AND source.ip:"[host]"}, sorted by \texttt{network.bytes} descending.
    \item \textit{Other hosts:} \texttt{event.module:suricata AND rule.name:"[signature]"}, all sensors.
    \item \textit{DNS tunneling suspicion:} \texttt{event.dataset:dns} sorted by query length or entropy, looking for unusually long subdomains.
    \item \textit{Cleartext creds:} \texttt{event.dataset:http AND http.request.method:POST}, review bodies for login forms over plain HTTP.
    \item Save a confirmed query as a saved search.
\end{itemize}

\section{HUNT: RED TEAM TRADECRAFT}
\begin{itemize}
    \item \textit{Automated push (Ansible or similar):} a burst of SSH (22/tcp) or WinRM (5985/5986/tcp) from one source to many hosts within seconds is a playbook run. \texttt{event.dataset:conn AND} \newline \texttt{destination.port:(22 OR 5985 OR 5986)}, grouped by source, sorted by distinct destination count.
    \item \textit{Impacket-style tooling:} short SMB sessions opening the \texttt{svcctl} or \texttt{winreg} named pipe then closing within seconds, across multiple hosts, one source. Needs Zeek's DCE-RPC analyzer: \texttt{event.dataset:dce\_rpc AND} \newline \texttt{rpc.endpoint:(svcctl OR winreg)}.
    \item \textit{Fan-out to admin shares:} SMB connections to \texttt{C\$}/\texttt{ADMIN\$} from a source that is not a known management host.
    \item A source hitting both the fan-out query and the admin-share query narrows red team's launch point fast.
\end{itemize}

\section{DATA SOURCES}
\begin{itemize}
    \item \textit{Suricata:} signature-based alerts.
    \item \textit{Zeek conn.log:} every connection seen, matched or not. Where movement/exfil show up first.
    \item \textit{Zeek dns.log:} every DNS query and response; good for tunneling and beaconing.
    \item \textit{Zeek http.log / ssl.log:} request/response metadata and certificate details, even without full payload capture.
    \item \textit{Zeek files.log:} files extracted from traffic, useful for catching a dropped payload in transit.
    \item \textit{Forwarded host logs:} account/process detail the network alone cannot see.
\end{itemize}

\section{WAZUH, IF DEPLOYED}
\begin{itemize}
    \item File Integrity Monitoring: fastest confirmation a specific file changed.
    \item Vulnerability detection: known-CVE software already installed.
    \item Scope query: \texttt{rule.groups:"syscheck" AND agent.name:"[host]"}.
    \item Custom rule example, extending a built-in FIM rule:
\end{itemize}
\begin{quote}
\footnotesize
\texttt{<rule id="100010" level="7">} \newline
\texttt{~~<if\_sid>550</if\_sid>} \newline
\texttt{~~<field name="file">nginx.conf</field>} \newline
\texttt{~~<description>config} \newline
\texttt{~~modified</description>} \newline
\texttt{</rule>}
\normalsize
\end{quote}

\section{WRITING A SURICATA RULE}
\begin{quote}
\texttt{alert tcp any any -> \$HOME\_NET 4444 (msg:"Possible} \newline \texttt{reverse shell"; sid:1000001; rev:1;)}
\end{quote}
Test against known-good and known-bad traffic before enabling live. A rule too broad fires constantly; too narrow misses variants.

\section{CLI HELPERS}
\begin{itemize}
    \item \texttt{zeek-cut}: pull specific fields out of raw Zeek TSV logs on the command line.
    \item \texttt{jq}: parse and filter JSON-formatted logs quickly if exporting from the SIEM.
    \item \texttt{tcpdump -nr [pcap] 'host [ip]'}: spot-check a specific host's traffic straight from a capture, no SIEM needed.
\end{itemize}

\section{ALERT TUNING}
\begin{itemize}
    \item Suppress or threshold a rule that fires on expected, benign traffic rather than disabling it outright.
    \item Keep a short list of intentionally suppressed rules, with the reason, so nobody re-enables noise blind.
    \item Alert fatigue: tune, do not just silence. A dropped real alert costs more than a noisy one.
\end{itemize}

\section{WAZUH ACTIVE RESPONSE}
\begin{itemize}
    \item Active response scripts can auto-react to a rule firing (e.g., blocking an IP), configured in \texttt{ossec.conf} on the manager.
    \item Start with logging-only responses during a competition; an automated action that misfires can do more damage than the thing it was meant to stop.
    \item \texttt{/var/ossec/logs/active-responses.log}: audit what active response has actually done.
\end{itemize}

\section{WRITING A WAZUH DECODER}
\begin{itemize}
    \item A decoder tells Wazuh how to parse a log line into fields a rule can then match against.
    \item Needed when a custom or unusual log format is not recognized out of the box.
    \item Test with \texttt{/var/ossec/bin/wazuh-logtest}, feeding it a sample line before trusting the decoder live.
\end{itemize}

\section{BUILDING A DASHBOARD PANEL}
\begin{itemize}
    \item Run and save the search first, then build a visualization from it (bar chart of matches over time is usually the most useful default).
    \item Add the visualization to a dashboard, and set a sensible default time range: last 24 hours for a daily check, last hour during an active incident.
    \item Name panels descriptively; "Panel 3" is useless to a teammate six hours into a round.
\end{itemize}

\section{LOG RETENTION}
\begin{itemize}
    \item Full packet capture fills disk fastest; decide the retention window before the round, not when the disk is already full.
    \item Elasticsearch index lifecycle settings (or manual index deletion) control how long searchable data sticks around.
    \item If disk fills mid-round, older indices are usually safer to drop than losing new ingestion; confirm this policy with the team ahead of time.
\end{itemize}

\section{CROSS-SOURCE CORRELATION}
\begin{itemize}
    \item A single alert is a clue. The same timestamp appearing in Suricata, Zeek, and a host log together is a much stronger case.
    \item When reporting scope to the Captain, cite which sources agree, not just the first alert that fired.
    \item Keep a short list of IOCs (IPs, hashes, filenames) confirmed bad this round; check new alerts against it fast.
\end{itemize}

\section{COMMON PITFALLS}
\begin{itemize}
    \item Storage vs. retention: decide what to keep before the disk fills.
    \item Agent drift: a disconnected Wazuh agent stops reporting silently; check the Agents page, not just alert volume.
    \item Treating "no alerts" as "nothing happening." Confirm sensors are actually capturing, not just quiet.
\end{itemize}

\section{PATTERN LIBRARY}
\begin{itemize}
    \item A sudden spike in DNS queries to one domain, especially with long or random-looking subdomains.
    \item Beaconing: a connection recurring at a suspiciously regular interval, to the same destination.
    \item A host generating alerts it has never generated before, even low-severity ones.
    \item An agent that stops checking in right around the same time a host starts behaving oddly elsewhere.
\end{itemize}

\end{multicols*}
\end{document}